\documentclass{article}
\usepackage[backend=biber, sorting=none]{biblatex}
\usepackage{graphicx}
\addbibresource{uni.bib}

\author{Angelo Battaglia}
\title{Proposta Tesi, tirocinio Argo Software}
\begin{document}

\maketitle

\begin{center}
\includegraphics[width=3in]{polito.jpg}
\includegraphics[width=3in]{og_logoargosoft.png}
\end{center}

\section{Introduzione}

Citando \textcite{StanfordOptimization}, un problema di ottimizzazione riguarda l'allocazione
delle risorse e l'analisi della complessità di un problema. L'ottimizzazione si suddivide in tre
sottoinsiemi: Linear Programming, Unconstrained Problems e Constrained Problems. In seguito, 
utilizzerò questi nomi sia in Italiano che in Inglese. Il successo delle tecniche di ottimizzazione
della programmazione lineare risiedono soprattutto nella formulazione del problema, piuttosto che
nel metodo utilizzato per risolverlo. Un problema con un grande numero di vincoli, infatti, 
presenta usualmente delle caratteristiche di linerità.

\subsection{Linear Programming}
Siano date \textit{i} entità, alle quali associamo \textit{\[w_i\]} pesi. Un problema base che riguarda 
la programmazione lineare è: \[massimizzare \quad  w_1x_1 + w_2x_2 \]
\[soggetto \ al \ vincolo \quad x_1 + x_2 \leq B \]
\[con \quad x_1\geq 0, \ x_2 \geq 0 \]

\subsection{Unconstrained Problems}
Problemi senza vincoli, ovvero unconstrained problems, sono quelli che considerano tutto il fascio 
di decisioni possibili. L'assenza di vincoli amplia la portata del problema, lo generalizza, aumentandone
i gradi di libertà. Inoltre, aggiungere dei vincoli a questa classe di problemi è un'operazione 
che non presenza, di solito, particolari criticità. A livello teorico, i problemi senza vincoli
assumono un particolare valore teorico, che successivamente va a concretizzarsi nella quotidianità
dei problemi con vincoli.

\subsection{Constrained Problems}
I problemi ai vincoli, detti Constrained Problems, trovano più utilità nelle applicazioni. I vincoli
vengono imposti nei problemi complessi per semplificarli, e per discretizzarli. Ogni vincolo 
riduce la portata del problema. Un problema generale che riguarda la programmazione matematica è: 
\[minimizzare \quad  f(\textbf{x}) \] 
\[soggetto \ ai \ vincoli \quad h_i(\textbf{x})=0, \ i= 1,2, ..., \ m \] \[g_j(\textbf{x})=0,\ \ j= 1,2, ..., \ p  \] 
\[con \  x \in \textit{S}\] 
In questa definizione, \textbf{x} è un vettore \textit{n}-dimensionale nelle incognite \(x_i\),
della forma \textbf{x} = \( (x_1, x_2, ..., x_n) \), ed \[ \textit{f} \ , \ \textit{\(h_i\)},  \ i= 1,2, ...,\ m,  
\ \textit{\(g_j\)}, \ j= 1,2, ..., \ p \] sono delle funzioni a valori reali delle variabili 
\(x_1, \ x_2, \ ..., \ x_n\).
Una misura ovvia della complessità di un problema è la sua dimensione. 
La dimensione di un problema è misurata nei termini del numero di incognite e
di vincoli. Come ci si può aspettare, i problemi trovano una più facile soluzione
quando la potenza dei calcolatori aumenta e la teoria matematica viene raffinata.
Un problema di larga scala viene ad avere migliaglia o milioni di variabili e vincoli.
Un possibile strada da seguire, in termini di software, è \textcite{Gecode}.

\subsection{Algoritmi Iterativi e convergenza}

La caratteristica più importante di un super-computer è la sua abilità di
computare operazioni ripetitive in maniera efficiente. Per questa ragione,
molti molti algoritmi di ottimizzazione, per loro natura, sono iterativi.
Tipicamente si parte da un vettore iniziale \(x_0\) per poi ottenerne uno, tramite un algoritmo, \(x_1\).
Il processo è ripetuto finché una soluzione migliore \(x_2\) viene trovata. Continuando finché 
Una sequenza di punti sempre migliori è trovata \( (x_0 , \ x_1 , ..., \ x_k , ...  )\) che si avvicini 
alla soluzione \textit{x*}. Per i problemi di programazione lineare risolti dal metodo del simplesso,
la sequenza generata è di lunghezza finita, e raggiunge la soluzione dopo un finito
numero di passi, anche se non inizialmente specificato.
Per la programmazione non lineare, la sequenza non raggiunge il punto \textit{x*} soluzione, 
ma converge verso esso, finché il processo termina ad un punto sufficientemente vicino al 
punto della soluzione.

\subsection{Convergence-rate theory}
Abbiamo due aspetti della convergence-rate theory. Il primo è generalmente 
conosciuto come analisi della complessità, e concerne la velocità dell'algoritmo 
di convergenza, distinguendo tra tempo polinomiale e tempo non polinomiale.
Il secondo aspetto provvede l'analisi di quanto il metodo velocemente converge 
nei suoi passi finali, e può provvedere un paragone dei vari metodi.

\section{Metodo del Simplesso}


\section{Tabu Search}

\subsection{Simulated Annealing}

Come cita l'articolo di  \textcite{SimulatedAnnealing01}, il Simulated Annealing
è un algoritmo di ottimizzazione che usa diversi strumenti matematici, come  

\section{Approccio Greedy}

\section{Algoritmi Genetici}

\section{Bin Packing}

\printbibliography[title = {Bibliografia e crediti}]

\end{document}
